\paragraph{So sánh LDA và QDA}

Linear Discriminant Analysis (LDA) và Quadratic Discriminant Analysis (QDA)
đều là các phương pháp phân loại dựa trên mô hình xác suất, với giả định rằng
dữ liệu của mỗi lớp tuân theo phân phối Gaussian đa biến. Điểm khác biệt cốt
lõi giữa hai phương pháp nằm ở giả định về ma trận hiệp phương sai, dẫn đến sự
khác biệt về dạng biên quyết định, số lượng tham số và khả năng tổng quát hóa.

\subparagraph{So sánh tổng quan}

\tablename~\ref{tab:lda-qda-comparison} trình bày các đặc điểm chính của hai
phương pháp.

\begin{table}[htbp]
  \centering
  \caption{So sánh giữa LDA và QDA}
  \label{tab:lda-qda-comparison}
  \import{\tablespath}{lda-qda-comparison}
\end{table}

\subparagraph{Phân tích theo số mẫu và số chiều}
Hiệu quả của LDA và QDA phụ thuộc mạnh vào mối quan hệ giữa số mẫu $n$ và số
chiều $d$ của dữ liệu:

\begin{itemize}
  \item \textbf{Trường hợp $n \gg d$:}
  \begin{itemize}
    \item Khi số mẫu lớn hơn nhiều so với số chiều, việc ước lượng các ma trận
    hiệp phương sai riêng $\boldsymbol{\Sigma}_k$ trong QDA trở nên chính xác.
    \item Khi đó, QDA có thể tận dụng sự linh hoạt của mình để mô hình hóa tốt
    các biên quyết định phi tuyến, thường cho hiệu năng cao hơn LDA.
  \end{itemize}

  \item \textbf{Trường hợp $n$ nhỏ hoặc $n \approx d$:}
  \begin{itemize}
    \item Việc ước lượng $\boldsymbol{\Sigma}_k$ trong QDA trở nên không ổn
    định và dễ dẫn đến ma trận suy biến.
    \item LDA có lợi thế vì chỉ cần ước lượng một ma trận hiệp phương sai
    chung $\boldsymbol{\Sigma}$, do đó ít tham số hơn và ổn định hơn.
    \item Trong trường hợp này, LDA thường tổng quát hóa tốt hơn.
  \end{itemize}

  \item \textbf{Trường hợp $d \gg n$:}
  \begin{itemize}
    \item Cả LDA và QDA đều gặp khó khăn trong việc ước lượng ma trận hiệp
    phương sai.
    \item Tuy nhiên, QDA bị ảnh hưởng nghiêm trọng hơn do cần ước lượng nhiều
    ma trận $\boldsymbol{\Sigma}_k$.
    \item Vì vậy, LDA thường là lựa chọn phù hợp hơn.
  \end{itemize}
\end{itemize}

\subparagraph{Kết luận}
Tóm lại, việc lựa chọn giữa LDA và QDA phụ thuộc vào đặc điểm của dữ liệu:
\begin{itemize}
  \item LDA phù hợp khi dữ liệu nhỏ, số chiều cao hoặc các lớp có phương sai
  tương tự nhau.
  \item QDA phù hợp khi dữ liệu đủ lớn và tồn tại sự khác biệt rõ rệt giữa
  phân phối của các lớp.
\end{itemize}

Sự lựa chọn này phản ánh sự đánh đổi giữa độ phức tạp của mô hình và khả năng
tổng quát hóa.
